%% kapa: exam=1 round=31 year=2020 period=2 form=A subject=appraisal_law subject_ko=감정평가관계법규 seq=1 number=1
%% source: Q-Net 감정평가사 기출문제 (cst00309) artlSeq=5207219 file=extracted/1008/5207219/A형/A형/2020년도 제31회 감정평가사 1차 시험_2교시_A형.pdf
%% conversion: kordoc 4.24.51 -> tools/gen1.py (auto); review=proofread-sonnet; flags=none
\begin{question}{1}
국토의 계획 및 이용에 관한 법령상 도시ㆍ군관리계획에 관한 설명으로 옳은 것은?
\begin{choices}
  \item 도시ㆍ군관리계획 결정의 효력은 지형도면을 고시한 날의 다음날부터 발생한다.
  \item 시ㆍ도지사는 국토교통부장관이 입안하여 결정한 도시ㆍ군관리계획을 변경하려면 미리 환경부장관과 협의하여야 한다.
  \item 도시ㆍ군관리계획을 입안할 수 있는 자가 입안을 제안받은 경우 그 처리 결과를 제안자에게 알려야 한다.
  \item 도시ㆍ군관리계획도서 및 계획설명서의 작성기준ㆍ작성방법 등은 조례로 정한다.
  \item 도지사가 도시ㆍ군관리계획을 직접 입안하는 경우 지형도면을 작성할 수 없다.
\end{choices}
\end{question}
\answer{3}
